%% 毕业设计小结

\begin{designsummary}
至此，我的本科毕业设计也接近尾声。回望这半年的过程，从最初确定题目、阅读文献，到后续完成实验、整理结果和撰写论文，我经历了一次较为完整的科研训练。作为本科阶段的最后一项重要任务，毕业设计不仅检验了我过去几年所学的专业知识，也让我更加真切地体会到，科研工作往往需要在不断尝试、反复修改和持续思考中逐渐推进。

进入实验室以来，在老师和同学们的帮助下，我逐渐接触到真正的科研工作。起初面对大量文献和不熟悉的问题时，常常会感到无从下手；在方法实现和实验验证过程中，也遇到过许多预想之外的困难。但正是在这些具体的问题中，我学会了如何梳理已有工作、如何抽象问题、如何根据实验结果反思方法，也更加明白了一个想法从提出到真正落地之间所需要经历的过程。

对我来说，毕业设计并不只是一篇论文，也像是本科四年学习和成长的一次总结。大学期间，我曾在迷茫中寻找方向，也曾在不同选择之间犹豫和试探。幸运的是，在这个过程中，我遇到了许多优秀的老师、同学和朋友，也在一次次交流、合作和自我调整中逐渐明确了自己想继续前进的道路。毕业设计凝结了这段时间里的努力，也记录了我从本科走向下一阶段之前的一次重要尝试。

当然，由于时间、经验和能力有限，这项工作仍然有许多不够完善的地方。无论是对问题的理解、方法的设计，还是论文的表达，都还有继续改进的空间。但也正因为如此，它更像是一个新的起点。未来我还将继续读博，面对更复杂、更长期的研究问题。希望自己能够记住完成这篇论文时的投入与不易，在之后的道路上保持好奇、耐心和继续探索的勇气。

本科生活即将结束，毕业设计也为这段时光留下了一个具体的注脚。它记录了我对科研的初步尝试，也记录了我在大学阶段的选择、成长与告别。未来再回望这段经历时，希望我依然会感谢此刻认真完成这份答卷的自己。
\end{designsummary}
